\scanpage{280}
\mspage{21}

Par \unsure{exemple|exple}, \unsure{prenant|prenons} le \deleted{chemin} type $\mathbb{\Delta}_2^+$ (i.e. $+1,+1$)
\diagram{\begin{tikzcd}[ampersand replacement=\&] 0 \arrow[r] \& 1 \arrow[r] \& 2 \end{tikzcd}},
correspondant aux chemins
\diagram{\begin{tikzcd}[ampersand replacement=\&] a_0 \arrow[r,dash,"{u_1}"] \& a_1 \arrow[r,"{u_2}"] \& a_2 \end{tikzcd}}
on peut associer à un tel chemin le chemin de type
\diagram{\begin{tikzcd}[ampersand replacement=\&] a_0 \arrow[r,"{u_2u_1}"] \& a_2 \end{tikzcd}},
ou aussi le chemin de type $\mathbb{\Delta}_2^+$
\diagram{\begin{tikzcd}[ampersand replacement=\&] a_0 \arrow[r,dash,"{u_2u_1}"] \& a_2 \arrow[r,"{\mathrm{id}}"] \& a_2 \end{tikzcd}},
\deleted{ou aussi} \deleted{\diagram{\begin{tikzcd}[ampersand replacement=\&] \unsure{a_0} \arrow[r,"{\mathrm{id}}"] \& a_0 \arrow[r,dash] \& \unsure{a_1} \arrow[r] \& {} \end{tikzcd}}} \deleted{\unsure{ou}}
etc. Comme le foncteur
\[ \left\{\begin{aligned}
X&\longmapsto \deleted{\unsure{\underline{\mathrm{Ch}}}}\mathrm{Ch}_\tau(X)\quad\unsure{(}(=\mathrm{Ob}\,\underline{\mathrm{Ch}}_\tau(X))\\
\mathrm{Cat}&\longrightarrow(\mathrm{Ens})
\end{aligned}\right. \]
est \uline{représenté} par l'objet $I_\tau$ de $X$, \unsure{le sorite|les sorites} des foncteurs représentables nous dit que de telles lois de formation \inserted{de chemins de type $\tau$ \deleted{\unsure{à}} à partir de chemins de type $\tau'$} correspondent \deleted{\unsure{biunivoquement}} \inserted{biunivoquement} \deleted{(à isom. unique près)} aux \deleted{\unsure{fonct}} flèches de $\mathrm{Cat}$
\diagram{\begin{tikzcd}[ampersand replacement=\&] I_{\tau'} \arrow[r,"{\varphi}"] \& I_\tau \end{tikzcd}},
en prenant simplement la ``loi'' correspondante à $\varphi$
\[ c\longmapsto c\circ\varphi\qquad\text{pour }c\colon I_\tau\longrightarrow X, \]
De façon \inserted{plus} précise, \deleted{\unsure{la correspondance}} \inserted{l'application} \deleted{\unsure{la loi}} naturelle
\[ \begin{aligned}
\mathrm{Hom}_{\mathrm{Cat}}(I_{\tau'},I_\tau)&\longrightarrow\mathrm{Hom}(\mathrm{Ch}_{\unsure{\tau}},\mathrm{Ch}_{\tau'})\\
\mathrm{Ob}\,\underline{\mathrm{Hom}}(I_{\tau'},I_\tau)&
\end{aligned} \]
% archive sheet number: 279 (last digit unclear)
